\label{ch:vacuum}

\section{Dérivation de l'échelle d'action à partir de l'état dodécaphasique}

La décennie d'action est déterminée intrinsèquement, sans l'adopter comme
unité héritée. La fonctionnelle déterminantale locale s'applique à la sortie HMT
préalablement construite \(\alpha_{\HMT}\) et définit
\begin{equation}
\Sigma_H=\varphi_{\HMT}\alpha_{\HMT}^{16}.
\label{eq:action-determinantal-reader}
\end{equation}
L'exposant \(16\) est l'ordre du conoyau local de la clôture signée ; sa valuation décimale fixe
\begin{equation}
\operatorname{ord}_{10}(\Sigma_H)=-34.
\label{eq:action-decade}
\end{equation}

L'origine de l'exposant peut être vérifiée dans ce volume. Sur
l'orbite locale à quatre feuillets agit le bloc de Hadamard
\begin{equation}
H_4=
\begin{pmatrix}
1&1&1&1\\
1&1&-1&-1\\
1&-1&1&-1\\
1&-1&-1&1
\end{pmatrix}.
\label{eq:action-hadamard}
\end{equation}

\begin{lemma}[Indice local de la droite d'action]
La forme normale de Smith de \(H_4\) est
\begin{equation}
\operatorname{SNF}(H_4)=\operatorname{diag}(1,2,2,4).
\label{eq:action-smith}
\end{equation}
En conséquence,
\[
\operatorname{coker}(H_4)\simeq
\Z/2\Z\oplus\Z/2\Z\oplus\Z/4\Z,
\qquad
|\operatorname{coker}(H_4)|=16.
\]
\end{lemma}

\begin{proof}
Le plus grand commun diviseur des entrées est un ; celui des mineurs
\(2\times2\) non nuls est deux ; celui des mineurs \(3\times3\) est quatre ; et
\(|\det H_4|=16\). Les quotients successifs de ces diviseurs déterminantaux
sont \(1,2,2,4\), qui prouvent \cref{eq:action-smith}. Le produit des
invariants donne l'ordre du conoyau.
\end{proof}

La norme de la section \(\alpha_{\HMT}\) sur ces seize feuillets
produit \(\alpha_{\HMT}^{16}\) ; l'auto-échelle \(\varphi_{\HMT}\) agit
une seule fois sur la base, et non une fois par feuillet. D'où
\cref{eq:action-determinantal-reader}. Les comparaisons certifiées
\begin{equation}
\alpha_{\HMT}^{16}<10^{-34}
<\varphi_{\HMT}\alpha_{\HMT}^{16}<10^{-33}
\label{eq:action-decade-bounds}
\end{equation}
démontrent, sans introduire l'unité SI, que l'ordre décimal est
exactement \(-34\). L'évaluation terminale commence par
\begin{equation}
\Sigma_H
=1.046243838104756580184229208303089\ldots\times10^{-34}.
\label{eq:action-sigma-value}
\end{equation}
Le remplacement du conoyau par trois copies avec exposant \(16^3\), ou
l'application réitérée de \(\varphi\) dans chaque feuillet, modifierait la
généalogie. Les deux opérations constituent des critères de réfutation et non
des conventions équivalentes.
La chaîne causale qui est conservée est donc
\[
\APP\to\TRIT\to\TPK\to U_{12}^{\rm sign}
\to\alpha_{\HMT}\to\Sigma_H
\to10^{-34}\U_S\to
\{\hbar_{\rm pre},\hbar_{\rm ret}\}.
\]
Ce volume ne redérive pas \(\alpha_{\HMT}\), mais développe la géométrie
métrologique postérieure à la sélection du type et de la décennie par
\(\Sigma_H\). L'entier \(-34\) n'est ni une mantisse ni un ajustement au Système
international, mais la valuation terminale de la fonctionnelle locale.

\section{Sections orientées de la droite d'action}

Les réalisations antérieure et postérieure au retour constituent deux sections
d'une même droite d'action :

\begin{equation}
\hbar_{\rm pre}=H_5\,10^{-34}\U_S,\qquad
\hbar_{\rm ret}=\eta_{\rm ret}\,10^{-34}\U_S.
\label{eq:action-twoway}
\end{equation}

La décennie commune détermine la droite, tandis que l'information
orientée s'exprime par

\begin{equation}
R_{\rm act}=\frac{\hbar_{\rm pre}}{\hbar_{\rm ret}}
=\frac{H_5}{\eta_{\rm ret}},
\qquad
\xi_{\rm act}=\frac12\log R_{\rm act}.
\label{eq:action-rapidity}
\end{equation}

La première quantité fournit une coordonnée multiplicative et la seconde
une coordonnée additive de type rapidité. Elles ne représentent pas deux constantes de
Planck indépendantes, mais les deux orientations d'une même section
par rapport au retour.

Si \(h_a=2\pi\hbar_a\), \(a\in\{\mathrm{pre},\mathrm{ret}\}\), une grandeur proportionnelle à \(h_a^{1/2}\) varie comme \(R_{\rm act}^{1/2}\), une grandeur proportionnelle à \(h_a^{-1/2}\) varie comme \(R_{\rm act}^{-1/2}\), et un quotient où l'action s'annule est un invariant bidirectionnel.

